\section{詳細解説}

\subsection{閉じた形の $A(q)$, $B(q)$ と数値計算の一致}

付録~\ref{app:twobody} で導いた胴体＋手足 1 本の閉じた形
\[
A(q) = I_b + I_l + \mu\left(|\bm a|^2+\ell_c^2+2\ell_c|\bm a|\cos(q-\varphi_a)\right),\quad
B(q) = I_l + \mu\left(\ell_c^2+\ell_c|\bm a|\cos(q-\varphi_a)\right)
\]
を、Sim01 の人間大モデル（腕のみ）で数値計算と比べた（図~\ref{fig:D1_closed_form.png}）。
両者の差は最大 $1.8\times10^{-15}$ で、浮動小数点の丸め誤差の範囲で一致した。
$A$ は 5.78〜7.61 kg\,m$^2$、$\Gamma=B/A$ は 0.028〜0.142 の範囲で変わる。

\DFIG{D1_closed_form.png}{左：腕 1 本の $A(q)$, $B(q)$（線：閉じた形、点：数値計算）。右：接続 $\Gamma(q)=B/A$。腕を真上に上げたとき（$q=90^\circ$）に最大、真下に下ろしたとき（$q=-90^\circ$）に最小になる。}

\begin{example}[$\Gamma$ を手計算する]
腕を真下に下ろした $q=-90^\circ$ では、肩が胴体重心の真上（$\varphi_a=90^\circ$）にあるので $\cos(q-\varphi_a)=\cos(-180^\circ)=-1$。
\begin{align*}
A &= 5.613 + 0.163 + 3.75\,(0.1225+0.1225-0.245) = 5.776,\\
B &= 0.163 + 3.75\,(0.1225-0.1225) = 0.163 ,
\end{align*}
$\Gamma = 0.0283$。腕を真上に上げた $q=90^\circ$ では $\cos=1$ なので $A=7.613$、$B=1.082$、$\Gamma=0.142$ である。
腕を下ろしているときは、腕の重心が肩の真下で胴体重心に近いため、腕を振っても胴体はあまり回らない。
腕を上げているときは、腕の重心が胴体重心から遠く、同じ角度だけ振っても胴体が約 5 倍回る。
\end{example}

\subsection{$A$ と $B$ の物理的な意味}

\begin{itemize}
\item $A(q)$ は「関節を固定して全体を剛体として回したときの慣性モーメント」である。腕を伸ばして胴体重心から遠ざけるほど大きい。
\item $B(q)$ は「胴体を止めたまま関節だけを回したときに生じる角運動量」である。腕の重心が胴体重心から遠い位置で回るほど大きい。
\item $\Gamma = B/A$ は「関節を 1 rad 回したときに胴体が逆向きに回る角度」で、0 と 1 の間の値をとる。
  手足が胴体に比べて十分軽ければ $\Gamma\to0$（胴体はほとんど回らない）、手足が非常に重ければ $\Gamma\to1$（胴体と手足が逆向きに同じだけ回る）に近づく。
\end{itemize}

\subsection{時間に依存しないことの意味}

基本式 $\dot\theta=-\Gamma(q)\dot q$ の両辺に $dt$ を掛けると $d\theta=-\Gamma(q)\,dq$ となり、時間が消える。
したがって、
\begin{itemize}
\item 同じ経路をゆっくり動かしても速く動かしても、胴体の回転角は同じである。
\item 速さが変えるのは、必要なトルク（$\propto f^2$、第~\ref{ch:dyn}章）と、かかる時間だけである。
\end{itemize}
この性質は、AMBAC の計画を「経路（形状空間の幾何）」と「時間配分（動力学）」に分けて考えてよいことを意味する。
第~\ref{ch:ai}章で、形状のループの形を学習しつつ、速さを別のフィードバック則で決める方策（軌道生成器型）が成功したのは、この分離と整合している。

\subsection{3 つの「できないこと」}

角運動量保存から、AMBAC には原理的にできないことが 3 つある。
\begin{enumerate}
\item \textbf{全角運動量を変えること}：回転を始めること・止めることは、外力なしにはできない（第~\ref{ch:rw}章の飽和）。
\item \textbf{重心を動かすこと}：手足をどう動かしても、系全体の重心の運動は変わらない。
\item \textbf{手足 1 本（平面）で向きを累積させること}：関節を戻せば向きも戻る（ホロノミック）。
\end{enumerate}
逆に、これら以外の「向きを変える」ことは、手足が 2 本以上あれば（3 次元なら 2 自由度の肩 1 つでも）任意の角度まで可能である。

\section*{考察課題}
\begin{exercise}
付録~\ref{app:twobody} の式で、肩を胴体重心に一致させる（$\bm a=0$）と、$\Gamma$ が $q$ によらず一定になる。このとき AMBAC は何に似た装置になるか。第~\ref{ch:rw}章と比べよ。
\end{exercise}
\begin{exercise}
腕の質量を 2 倍にしたとき、$\Gamma$ の最大値はおよそ何倍になるか。換算質量 $\mu$ の変化に注意して見積もれ。
\end{exercise}
